

\chapter{Conclusão}
\label{cha:conclusao}

A metodologia de aplicação do GPR envolve uma série de etapas, desde o planejamento da pesquisa até a interpretação dos resultados. A correta execução dessas etapas é essencial para obter informações precisas e confiáveis sobre as características subsuperficiais. A combinação de conhecimento geológico, experiência prática e o uso adequado de técnicas de processamento de dados são fundamentais para o sucesso da aplicação do GPR.

Neste trabalho apresenta-se uma metodologia para inversão em eletrorresistividade usando o \ac{aog}. Os resultados mostraram que o algoritmo pode ser usado em inversões 1D, embora, com o aumento da quantidade camadas, é preciso aumentar a quantidade de agentes (gafanhotos) para se ter uma boa resolução do problema inverso. O tempo de execução do algoritmo para modelos com camadas acima de três pode ser custoso, dependendo da máquina que está sendo executada.  
O \ac{aog} pode ser usado o paralelismo para lidar com essas questões ou utilizar variações do algoritmo, ou versões híbridas, no trabalho de \citeonline{meraihi2021grasshopper} apresenta várias abordagens e usos do \ac{aog} em diversas áreas.